% ============================================================================
% CONCLUSION SECTION (Paper-Ready)
% ============================================================================
% Status: FINAL - Kurz, prägnant, keine neuen Claims
% Date: 2025-12-28
% Placement: After Discussion, before References
% ============================================================================

\section{Conclusion}

We have introduced a formal distinction between \textbf{latent knowledge} and \textbf{knowledge activation} in large language models. Our central finding is that models can possess knowledge---verified through direct retrieval---yet systematically fail to apply it during generative tasks. This gap persists even under maximally explicit epistemic constraints.

We formalize this observation through the concept of \textbf{activation competence} ($\alpha$), defined as the conditional probability of correct knowledge application given knowledge availability:
\[
\alpha := \mathbb{P}\big(A(x) \mid M_{\text{latent}}(x) = 1\big)
\]

The Activation Challenge provides a minimal, controlled diagnostic that isolates this dimension. Across four contemporary language models, we observe:
\begin{itemize}
    \item All models exhibit activation gaps under baseline conditions ($\alpha < 1$).
    \item Explicit epistemic prompting (ESL) increases $\alpha$ but does not guarantee $\alpha = 1$.
    \item Some models show bounded activation competence, indicating intrinsic limitations.
\end{itemize}

These results have implications for how we understand hallucination, evaluate model reliability, and design prompting strategies. They suggest that improving epistemic behavior may require interventions beyond prompt engineering---potentially at the level of training objectives or model architecture.

We hope this work contributes to a more principled understanding of when and why language models fail to use what they know, and provides a foundation for future research on knowledge activation in generative systems.


% ============================================================================
% OPTIONAL: Final sentence (can be used or omitted)
% ============================================================================
%
% \begin{quote}
% \textit{Knowing is not enough; the question is whether knowledge is activated.}
% \end{quote}
%
